\section*{Capítulo \ref{ch:g11:color-light-sources} --- El color y las fuentes de luz}

\begin{solution}{exo:g11:color-light-sources:1}
(a) amarillo; (b) cian; (c) blanco. Complementaria del azul: el amarillo
(azul $+$ amarillo $=$ azul $+$ rojo $+$ verde $=$ blanco).
\end{solution}

\begin{solution}{exo:g11:color-light-sources:2}
Amarillo: R, G. Magenta: R, B. Blanco: R, G, B. Negro: ninguno. Naranja:
R a plena \omterm{def:g10:energy-conservation:power}{potencia} y G a \omterm{def:g10:energy-conservation:power}{potencia} reducida.
\end{solution}

\begin{solution}{exo:g11:color-light-sources:3}
El magenta absorbe la banda verde: emerge $(R, B)$, luz magenta. El
\omterm{def:g11:color-light-sources:subtractive}{filtro} cian absorbe después la banda roja: solo sobrevive $(B)$ ---
azul.
\end{solution}

\begin{solution}{exo:g11:color-light-sources:4}
Franja amarilla: \omterm{def:g11:color-light-sources:objectcolor}{difunde} R y G, absorbe B. (a) Bajo luz roja \omterm{def:g11:color-light-sources:objectcolor}{difunde} R:
roja. (b) Bajo luz azul lo absorbe todo: negra.
\end{solution}

\begin{solution}{exo:g11:color-light-sources:5}
$T = \dfrac{\qty{2.90e-3}{m.K}}{\qty{5.80e-7}{m}} = \qty{5000}{K}$.
\end{solution}

\begin{solution}{exo:g11:color-light-sources:6}
\emph{(a)} $\lambda_{\max} = \qty{2.90e-3}{m.K}/\qty{1800}{K}
= \qty{1.61e-6}{m} = \qty{1610}{nm}$: \omterm{def:g10:light-spectra:wavelength}{infrarrojo}.

\emph{(b)} El \omterm{def:g10:light-spectra:white}{espectro} es continuo: su cola visible es pequeña pero no
nula, y el extremo azul es el primero que se apaga --- lo que queda
visible es sobre todo rojo anaranjado, de ahí el color de la llama.
\end{solution}

\begin{solution}{exo:g11:color-light-sources:7}
Chaqueta (\omterm{def:g11:color-light-sources:objectcolor}{difunde} R): (a) roja; (b) roja; (c) negra. Guantes (difunden G
y B, absorben R): (a) cian; (b) negros; (c) verdes.
\end{solution}

\begin{solution}{exo:g11:color-light-sources:8}
\emph{(a)} Un \omterm{def:g10:light-spectra:emission}{espectro de emisión de líneas}. Cualquier objeto solo puede
difundir $\qty{589}{nm}$ o nada: la calle es una gama de amarillos
anaranjados y negros.

\emph{(b)} El coche azul solo \omterm{def:g11:color-light-sources:objectcolor}{difunde} la banda azul. La luz del día la
contiene: coche azul. La línea del sodio no lleva nada de azul: todo se
absorbe y el coche se ve negro.
\end{solution}

\begin{solution}{exo:g11:color-light-sources:9}
\emph{(a)} El haz blanco se abre en un arcoíris; el haz del \omterm{def:g11:color-light-sources:lamps}{láser} se
desvía pero sigue siendo un único haz del mismo verde: es monocromático
--- una sola \omterm{def:g10:light-spectra:wavelength}{longitud de onda}, nada que el prisma pueda separar.

\emph{(b)} El \omterm{def:g11:color-light-sources:subtractive}{filtro} rojo transmite solo la banda roja y absorbe
$\qty{532}{nm}$: no sale casi nada; la mancha desaparece.
\end{solution}

\begin{solution}{exo:g11:color-light-sources:10}
\omterm{def:g10:light-spectra:white}{Espectro}: un pico azul estrecho cerca de $\qty{450}{nm}$ más la banda
amarilla ancha del \omterm{def:g11:color-light-sources:led}{fósforo} (aproximadamente de $500$ a
$\qty{700}{nm}$). El azul y el amarillo son complementarios, así que el
trío de los \omterm{def:g11:color-light-sources:cones}{conos} se lee como blanco. La banda se apaga hacia el rojo
profundo, de modo que un objeto rojo intenso recibe poco que pueda
difundir: se ve más apagado que a la luz del día, cuyo \omterm{def:g10:light-spectra:continuous}{espectro continuo}
lo alimenta por completo.
\end{solution}

\begin{solution}{exo:g11:color-light-sources:11}
(a) $\lambda_{\max} = \qty{2.90e-3}{m.K}/\qty{1800}{K}
= \qty{1610}{nm}$, \omterm{def:g10:light-spectra:wavelength}{infrarrojo};
(b) $\qty{2.90e-3}{m.K}/\qty{6500}{K} = \qty{446}{nm}$, azul violeta. La
luz de la vela es la que se llama ``cálida'' --- la fuente de
temperatura \emph{más baja}: el vocabulario va al revés que la escala
\omterm{def:g10:light-spectra:kelvin}{kelvin}.
\end{solution}

\begin{solution}{exo:g11:color-light-sources:12}
\emph{(a)} \omterm{def:g11:color-light-sources:perception}{Metámeros}. Su existencia demuestra que el color es una
\omterm{def:g11:color-light-sources:perception}{percepción} --- tres respuestas de los \omterm{def:g11:color-light-sources:cones}{conos} --- y no el \omterm{def:g10:light-spectra:white}{espectro} en sí:
infinitos \omterm{def:g10:light-spectra:white}{espectros} se proyectan sobre el mismo trío.

\emph{(b)} Prisma o \omterm{def:g10:light-spectra:white}{espectroscopio}: una línea en $\qty{589}{nm}$ frente
a dos bandas en $540$ y $\qty{610}{nm}$. \omterm{def:g11:color-light-sources:subtractive}{Filtro} rojo: absorbe
$\qty{589}{nm}$, la mancha del sodio se apaga, mientras que el subpíxel
de $\qty{610}{nm}$ de la pantalla sobrevive --- la zona de pantalla se
vuelve roja.
\end{solution}

\begin{solution}{exo:g11:color-light-sources:13}
\emph{(a)} $\lambda_{\max} = \qty{2.90e-3}{m.K}/\qty{3200}{K} =
\qty{906}{nm}$. \omterm{def:g11:color-light-sources:colortemp}{Temperatura de color}
$\qty{3200}{K} < \qty{6500}{K}$: más cálida (más anaranjada) que la luz
del día.

\emph{(b)} La gelatina absorbe parte del exceso rojo anaranjado para
reequilibrar el \omterm{def:g10:light-spectra:white}{espectro} hacia el azul; y como un \omterm{def:g11:color-light-sources:subtractive}{filtro} solo quita luz,
la \omterm{def:g10:energy-conservation:power}{potencia} total sobre el escenario baja.

\emph{(c)} Un \omterm{def:g11:color-light-sources:subtractive}{filtro} es pasivo: puede absorber en cada \omterm{def:g10:light-spectra:wavelength}{longitud de onda},
pero nunca emitir, así que ninguna gelatina puede subir la intensidad
azul por encima de la que el \omterm{def:g11:lenses-and-eye:focal}{foco} ya produce.
\end{solution}

\begin{solution}{exo:g11:color-light-sources:14}
\emph{(a)} $\qty{2700}{K}$: luz cálida, rica en rojo --- la banda roja
de la carne está bien alimentada y la carne se ve de un rojo vivo y
fresco. $\qty{6500}{K}$: blanco azulado y frío --- el pescado \omterm{def:g11:color-light-sources:objectcolor}{difunde} el
trío completo y se ve de un blanco brillante, recién puesto en hielo.

\emph{(b)} Ninguno de los dos ha cambiado su mercancía: las costumbres
de absorción de la carne y del pescado siguen intactas. Lo que ha
cambiado es la iluminación y, con ella, la luz difundida --- el \omterm{def:g11:color-light-sources:objectcolor}{color de
un objeto} solo está definido bajo una iluminación dada.
\end{solution}

\begin{solution}{exo:g11:color-light-sources:15}
\emph{(a)} El sodio: todo el resplandor de la ciudad está en una línea
de $\qty{589}{nm}$, que un único \omterm{def:g11:color-light-sources:subtractive}{filtro} estrecho retira de las imágenes
del telescopio. Un \omterm{def:g10:light-spectra:white}{espectro} led ancho no se puede filtrar sin tirar
también la luz de las estrellas.

\emph{(b)} Conductores y peatones prefieren el led: bajo su \omterm{def:g10:light-spectra:white}{espectro}
completo, los objetos conservan colores distinguibles (un abrigo rojo
sigue siendo rojo), mientras que la luz de sodio lo vuelve todo amarillo
anaranjado o negro.

\emph{(c)} Ledes blancos cálidos, en torno a $\qty{2700}{K}$: una
reproducción de color aceptable para la calle y un \omterm{def:g10:light-spectra:white}{espectro} pobre en el
extremo azul, que es la parte que más contamina el cielo nocturno.
\end{solution}

\begin{solution}{pb:g11:color-light-sources:1}
\textbf{1.} R $+$ G $=$ amarillo, R $+$ B $=$ magenta, G $+$ B $=$
cian; la superposición triple, blanco.

\textbf{2.} Naranja: rojo a plena \omterm{def:g10:energy-conservation:power}{potencia}, verde a \omterm{def:g10:energy-conservation:power}{potencia} parcial y
azul apagado. Rosa pálido: los tres encendidos (base blanca) con el rojo
ligeramente dominante.

\textbf{3.} La sombra que proyecta el \omterm{def:g11:lenses-and-eye:focal}{foco} rojo solo recibe el \omterm{def:g11:lenses-and-eye:focal}{foco}
verde: sombra verde; simétricamente, la otra sombra es roja.

\textbf{4.} La pintura azul oscura solo \omterm{def:g11:color-light-sources:objectcolor}{difunde} B. Bajo el \omterm{def:g11:lenses-and-eye:focal}{foco} rojo
solo: negra. Bajo los tres: \omterm{def:g11:color-light-sources:objectcolor}{difunde} la banda azul --- azul.

\textbf{5.} El ojo reduce cualquier \omterm{def:g10:light-spectra:white}{espectro} a tres respuestas de los
\omterm{def:g11:color-light-sources:cones}{conos}; basta con igualar el trío para igualar el color. Los tres tipos
de \omterm{def:g11:color-light-sources:cones}{conos} fijan el número de \omterm{def:g11:lenses-and-eye:focal}{focos}.

\textbf{6.} Amarillo $=$ blanco $-$ azul: el tinte absorbe la banda
azul.

\textbf{7.} \omterm{def:g11:lenses-and-eye:focal}{Foco} rojo: \omterm{def:g11:color-light-sources:objectcolor}{difunde} R --- rojo. \omterm{def:g11:lenses-and-eye:focal}{Foco} verde: \omterm{def:g11:color-light-sources:objectcolor}{difunde} G ---
verde. \omterm{def:g11:lenses-and-eye:focal}{Foco} azul: lo absorbe todo --- negro.

\textbf{8.} Iluminando la escena solo con el \omterm{def:g11:lenses-and-eye:focal}{foco} azul; el golpe de
interruptor es blanco $\to$ azul. El traje amarillo \omterm{def:g11:color-light-sources:objectcolor}{difunde} R y G pero
absorbe B: brillante bajo el blanco y negro bajo el azul.

\textbf{9.} Gelatina magenta: $(R, G, B) \to (R, B)$. El traje \omterm{def:g11:color-light-sources:objectcolor}{difunde} R
y G de lo que recibe: \omterm{def:g11:color-light-sources:objectcolor}{difunde} R y absorbe B --- traje rojo.

\textbf{10.} Una gelatina roja conserva una banda de tres y tira dos
tercios de la \omterm{def:g10:energy-conservation:power}{potencia} del \omterm{def:g11:lenses-and-eye:focal}{foco} blanco; cada gelatina C, M o Y quita una
sola banda, conserva dos tercios, y apilando dos de ellas se sigue
obteniendo cualquier primario.

\textbf{11.} Un \omterm{def:g10:light-spectra:emission}{espectro de emisión de líneas}. La calle se convierte en
un mundo monocromo: cada superficie es amarillo anaranjada, más apagada
o más viva, o negra.

\textbf{12.} Pared blanca: amarillo anaranjada (\omterm{def:g11:color-light-sources:objectcolor}{difunde} la línea). Marca
vial amarilla: amarillo anaranjada y viva. Coche rojo: oscuro --- la
línea de $\qty{589}{nm}$ queda fuera de la banda que \omterm{def:g11:color-light-sources:objectcolor}{difunde}. Coche
azul: negro.

\textbf{13.} Un único \omterm{def:g11:color-light-sources:subtractive}{filtro} estrecho centrado en $\qty{589}{nm}$ retira
de una imagen todo el resplandor de la ciudad sacrificando casi nada de
los \omterm{def:g10:light-spectra:white}{espectros} anchos de las estrellas.

\textbf{14.} $\lambda_{\max} = \qty{2.90e-3}{m.K}/\qty{2700}{K} =
\qty{1.07e-6}{m} = \qty{1070}{nm}$: el máximo, y casi toda la \omterm{def:g10:energy-conservation:power}{potencia},
están en el \omterm{def:g10:light-spectra:wavelength}{infrarrojo} --- calor, no luz.

\textbf{15.} Sodio: $0.30 \times \qty{100}{W} = \qty{30}{W}$ de luz
visible; wolframio: $0.05 \times \qty{100}{W} = \qty{5}{W}$. Seis veces
más luz por la misma factura.

\textbf{16.} El \omterm{def:g11:color-light-sources:led}{fósforo} convierte parte del azul en una banda amarilla
ancha; el azul y el amarillo son complementarios, así que azul $+$
amarillo completa el trío: blanco.

\textbf{17.} Un pico azul estrecho más una joroba amarilla ancha ---
frente al continuo uniforme de la luz del día. El extremo rojo profundo
(y un valle entre el azul y el amarillo) queda infrarrepresentado.

\textbf{18.} Los objetos rojos intensos se reproducen apagados y
parduzcos. De ahí los ledes cálidos de $\qty{2700}{K}$ del carnicero,
cuyo contenido rojo reforzado mantiene viva la carne.

\textbf{19.} Nada de su interior está a $\qty{2700}{K}$: el chip
funciona cerca de la temperatura ambiente. La etiqueta solo promete el
\emph{tono} de un cuerpo incandescente a $\qty{2700}{K}$, cuyo máximo
estaría en
$\qty{2.90e-3}{m.K}/\qty{2700}{K} = \qty{1.07e-6}{m} = \qty{1070}{nm}$.

\textbf{20.} Dormitorio: unos $\qty{2700}{K}$ --- luz cálida y pobre en
azul, para descansar. Joyero: de $5500$ a $\qty{6500}{K}$ --- un blanco
parecido a la luz del día, para que los blancos se lean blancos y las
piedras devuelvan todas las bandas. La ley: color percibido $=$ \omterm{def:g10:light-spectra:white}{espectro}
de la fuente $\times$ absorción del objeto, leído a través de tres
\omterm{def:g11:color-light-sources:cones}{conos}.
\end{solution}
